% Conclusions (body only; master supplies \section{Conclusions}).
% Distinct contribution: base-drag closure choice dominates integrated
% supersonic-rocket apogee error (8.10 pp), benchmarked + generalization-tested.
% Future work (prioritized): independent finned-base base-pressure dataset
% (B->A on FINNED_BASE_K), transonic blend fix, Sahu digitization.
% Every number per paper/_shared/CANONICAL_FACTS.md (sections B, C, D, K).
% No Devan-Ashwood: Cd_base = 0.064 + 0.186/M^2 is an empirical correlation
% anchored on NACA TN 3393 + ESDU 77021. \cite{yu2026jsr} is a placeholder key.
This study set out to answer a question distinct from the integrated accuracy
assessment of the companion flagship paper~\cite{yu2026jsr}: among the supersonic
submodels of an open-source rocket simulator, how heavily does the base-drag term
weigh on the integrated apogee prediction relative to the others, and can the
closure that carries that weight be anchored to external base-pressure data? By
re-running the flight corpus with each supersonic submodel disabled in isolation,
we established that the base-drag term is not a second-order detail but the
dominant driver of apogee error. Disabling the finned-base drag augmentation ---
the base-drag component the corpus is most sensitive to --- moves predicted apogee
by a mean of $8.10$~pp across 24 of the 25 corpus flights (MESOS~293K excluded;
$23$ single-stage plus the AeroPac~104K two-stage closure; peaking at $39.5$~pp on
the Mach~$2.19$ A-601 Kinsel), an order of magnitude larger than the next
supersonic submodel --- Van~Driest~II compressible skin friction at $0.87$~pp
\cite{hopkins1971} --- and far above the DATCOM~4.1.5.1 fin wave-drag term at
$0.39$~pp and the structurally prominent shock-geometry pre-pass at only
$0.15$~pp. The pre-pass
earns its place through local-flow fidelity for fin loads and stability, not
through any gross-apogee effect; the apogee budget is dominated instead by how
the closure treats the finned base. This reframes where modeling effort and
validation scrutiny should be concentrated in supersonic-rocket apogee
prediction.

Crucially, the closure branches benchmarked here are anchored on published
base-pressure data rather than tuned in a vacuum, which is what licenses the
central claim as physics rather than as an artifact of one calibration. The
empirical turbulent correlation $C_{d,\text{base}} = 0.064 + 0.186/M^2$
reproduces the NACA~TN~3393 nonlifting-body base-pressure measurements at a MAPE
of $15.9\%$ (and $4.0\%$ against the Hart NACA~RM~L52E06 free-flight
subset~\cite{hart1952basedrag}) and is consistent with the ESDU~77021
methodology \cite{chapman1955,esdu77021}; the Chapman laminar branch matches
the same data at $4.4\%$ \cite{chapman1950}. The Chapman--Korst turbulent
shear-layer treatment \cite{chapman1958}, the Viswanath boattail relief
\cite{viswanath1996}, and the DATCOM~4.1.5.1 fin wave-drag closure
\cite{datcom1978} complete the set against which the finned-base augmentation
is judged. The dominant closure is also the one carrying the most empirical
calibration, and we confront that head-on rather than asking for trust: in the
decontaminated prospective holdout --- in which every flight that touched a
scale constant (Raven, Rabia, Rabia Short~Fin~Can, Kinsel, Torrent) is forced
into the development split --- the genuinely blind holdout subset (MAE~$3.95\%$,
$n=12$) is \emph{more} accurate than the development subset (MAE~$5.47\%$,
$n=13$). A closure that improves out of sample is generalizing the underlying
base-flow physics, not memorizing the calibration flights; this inversion is
the primary evidence that neutralizes the in-sample concern attached to the
dominant mechanism.

These results carry a clear prioritization for future work. First, and most
important, the finned-base augmentation is the dominant apogee driver yet
remains a corpus-anchored, B-level closure: it is calibrated against the
integrated flight corpus rather than isolated against a dedicated component
dataset. An independent finned-base base-pressure dataset --- measuring base
pressure on a finned afterbody across the relevant supersonic Mach range,
ideally on a flat-base geometry comparable to the Basic Finner reference
\cite{dupuis1997} --- would promote this constant from B-level to A-level
external validation and is the single highest-value evidence gap for the
apogee question this paper raises. Second, the transonic blend remains the
disclosed weak band: it is the only regime carrying a one-sided mean signed
apogee bias ($-3.66\%$), and there the high-supersonic empirical correlation,
the Chapman--Korst shear-layer treatment, and the transonic peak polynomial
must be reconciled across the Mach~$0.8$--$1.3$ window. Tightening this blend
against a dedicated transonic base-pressure benchmark is the most direct route
to reducing the residual apogee error. Third, the published transonic
base-flow computation of Sahu, Nietubicz, and Steger \cite{sahu1983} on a
secant-ogive--cylinder--boattail geometry would furnish a third-party CFD
comparator for the boattail and transonic closures once its base-pressure
curves are digitized; that digitization is deferred here because the geometry
differs structurally from the finned reference body and warrants a dedicated
model build, but closing it is the natural complement to the transonic blend
work.

In sum, this open mechanism-attribution study establishes a result that the
integrated parity headline of the companion paper \cite{yu2026jsr} could not
surface on its own: for supersonic-rocket apogee, the base-drag term ---
specifically the treatment of the finned base --- is where accuracy is won or
lost, dominating the supersonic error budget by an order of magnitude over every
other submodel. The closure branches are benchmarked against externally measured
base-pressure data and the corpus-frozen constants are shown to generalize out of
sample, and the path to elevating the dominant closure from a corpus-anchored
constant to an externally validated submodel is now explicit. All code, calibration constants, the ablation harness, and the
underlying flight corpus are released openly (Section~9) so that the
intercomparison can be reproduced and the identified evidence gaps closed by
the community.
